% chapters/04/section-04.tex
\section{Siklus Lengkap}
\label{sec:siklus-lengkap}

Bagian ini menyatukan seluruh unsur yang telah dibahas menjadi program siklus
permainan yang lengkap. Susunan keempat unsurnya diperlihatkan pada
Gambar~\ref{fig:siklus-dengan-clock}.

\begin{figure}[htbp]
\centering
% chapters/04/figures/siklus-dengan-pengendali-waktu.tikz
% Siklus lengkap: tiga tahap pekerjaan ditambah pengendali waktu yang
% mengukur selang waktu dan menahan laju bingkai.
\begin{tikzpicture}[
  tahap/.style={
    rectangle, rounded corners=2pt, draw=ubPrimary, thick, fill=ubAlt,
    align=center, text width=25mm, inner sep=3.5pt, minimum height=18mm,
    font=\footnotesize
  },
  waktu/.style={tahap, draw=ubAccent, fill=white},
  panah/.style={-{Stealth[length=2.4mm]}, thick, draw=ubAccent}
]

\node[tahap] (input) at (0,0)
  {\textbf{1. Masukan}\\[2pt]\scriptsize peristiwa dari\\pemain};
\node[tahap, right=7mm of input] (update)
  {\textbf{2. Pembaruan}\\[2pt]\scriptsize keadaan\\permainan};
\node[tahap, right=7mm of update] (render)
  {\textbf{3. Penayangan}\\[2pt]\scriptsize gambar ulang\\seluruh jendela};
\node[waktu, right=7mm of render] (clock)
  {\textbf{4. Pengendali waktu}\\[2pt]\scriptsize ukur selang,\\tahan laju bingkai};

\draw[panah] (input) -- (update);
\draw[panah] (update) -- (render);
\draw[panah] (render) -- (clock);
\draw[panah] (clock.south) -- ++(0,-8mm) -| (input.south);

\node[font=\scriptsize\itshape, ubPrimary, anchor=north] at (4.8,-2.2)
  {diulang selama permainan berjalan};
\end{tikzpicture}

\caption{Siklus lengkap: tiga tahap pekerjaan dan pengendali waktu}
\label{fig:siklus-dengan-clock}
\end{figure}

\subsection{Kode Lengkap}

\kodepython{python/ch04_game_loop.py}{Siklus permainan lengkap beserta pengendali waktu}{lst:ch04-game-loop}

\subsection{Menemukan Ketiga Tahap di Dalam Kode}

Ketiga tahap siklus sengaja diberi keterangan pada kode di atas agar mudah
ditemukan. Perhatikan susunannya.

\begin{description}
  \item[Tahap 1: membaca masukan] Terdapat pada perulangan \code{for event in
        pygame.event.get()}. Perulangan itu mengosongkan antrean peristiwa dan
        menangani setiap peristiwa yang ada: menutup jendela, menekan ESC,
        mengubah batas laju dengan tombol panah, atau mengembalikan penghitung
        bingkai dengan tombol spasi.
  \item[Tahap 2: memperbarui keadaan] Terdapat pada pemanggilan
        \code{update_state}, penambahan penghitung bingkai, pembacaan laju
        terukur, dan pembaruan judul jendela. Pada bab ini keadaan yang
        diperbarui masih berupa penghitung bingkai; mulai Bab~V keadaan yang
        diperbarui mencakup posisi dan kecepatan objek.
  \item[Tahap 3: menayangkan] Terdapat pada pemanggilan \code{draw_window}
        dan \code{pygame.display.flip()}. Fungsi \code{draw_window} mengisi
        latar, menempelkan seluruh teks, dan menggambar batang penanda.
  \item[Pengendali waktu] Terletak di akhir putaran, yaitu pada
        \code{clock.tick(fps_limit)}. Perintah itu menahan sisa waktu satu
        bingkai apabila pekerjaan sudah selesai lebih cepat, dan mengembalikan
        selang waktu yang benar-benar berlangsung.
\end{description}

\subsection{Tiga Kebiasaan yang Perlu Dibangun}

Program tersebut memuat tiga kebiasaan yang berlaku pada seluruh bab
berikutnya.

\textbf{Pertama, gambar seluruh jendela dari awal setiap bingkai.} Perhatikan
bahwa \code{draw_window} selalu diawali \code{screen.fill()}, yang menimpa
seluruh isi jendela dengan satu warna. Apabila langkah itu dihilangkan,
gambar dari bingkai sebelumnya tetap tertinggal dan objek yang bergerak
meninggalkan jejak di sepanjang lintasannya. Cara berpikir yang harus
dibangun adalah: setiap bingkai dimulai dari halaman kosong.

\textbf{Kedua, simpan keadaan di luar perulangan.}
\code{frame_count}, \code{fps_limit}, dan \code{start_ticks} dibuat sebelum
perulangan dimulai. Ketiganya harus hidup dari satu bingkai ke bingkai
berikutnya. Apabila \code{frame_count = 0} dipindahkan ke dalam perulangan,
penghitung itu kembali ke nol setiap bingkai dan tidak pernah bertambah.

\textbf{Ketiga, ubah keadaan hanya pada tahap pembaruan.}
Penambahan \code{frame_count} hanya terjadi pada tahap kedua, bukan di dalam
fungsi penggambaran. Memisahkan ``mengubah keadaan'' dari ``menggambar
keadaan'' membuat program lebih mudah diperiksa: fungsi penggambaran tidak
mengubah apa pun, sehingga dapat dipanggil berapa kali pun tanpa
mempengaruhi jalannya permainan.

\begin{peringatanbox}
Kesalahan yang paling sering terjadi pada tahap ini adalah menaruh
perubahan keadaan di dalam tahap penayangan. Misalnya, menambahkan
penghitung bingkai di dalam fungsi penggambaran. Akibatnya penghitung itu
bertambah setiap kali gambar disusun, bukan setiap bingkai, dan jumlahnya
berbeda-beda menurut berapa kali fungsi tersebut dipanggil. Jaga batas
antara kedua tahap itu dengan disiplin.
\end{peringatanbox}

\subsection{Mengamati Program Bekerja}

Ketika program dijalankan, tiga hal dapat Anda amati untuk memastikan siklus
bekerja sebagaimana mestinya.

\begin{itemize}
  \item \textbf{Judul jendela berubah setiap bingkai}, menampilkan laju
        terukur pada saat itu. Judul yang tidak berubah menandakan siklus
        tidak berjalan.
  \item \textbf{Batang penanda bertambah panjang} lalu kembali ke nol setelah
        mencapai tepi kanan. Batang itu memperlihatkan bahwa tahap pembaruan
        dan penayangan benar-benar dijalankan berulang kali.
  \item \textbf{Laju terukur mendekati batas yang diminta}, misalnya 59,8
        ketika batasnya 60. Selisih kecil itu wajar, sebab pengukuran
        dilakukan atas waktu yang sesungguhnya berjalan.
\end{itemize}

Tekan tombol panah atas dan bawah untuk mengubah batas laju, lalu amati
perubahan angka pada judul jendela. Percobaan sederhana tersebut memperlihatkan
hubungan antara batas yang diminta dan laju yang tercapai.
